\begin{figure}[htbp]
  \centering
  \begin{adjustbox}{max width=\textwidth}
    \begin{tikzpicture}[
        % Distanze calibrate
        node distance=4.5ex and 2.5em,
        %
        % Stili di base
        base/.style={draw, rounded corners=2pt, inner sep=8pt, align=center, font=\small},
        %
        % Core invariante leggermente allargato a 22em per accogliere la nuova riga dei contratti
        core/.style={base, fill=blue!8, draw=blue!45, minimum width=22em, thick},
        % Plugin mantenuto proporzionale (2em in meno del core)
        active/.style={base, fill=green!10, draw=green!50!black, minimum width=20em},
        %
        % Moduli Satelliti
        planned/.style={base, fill=gray!4, draw=gray!40, dashed, minimum width=8em},
        %
        % Stili Frecce
        arr-active/.style={-latex, thick, green!50!black},
        arr-planned/.style={-latex, thick, dashed, gray!50}
      ]

      %% ── Core immutabile (Centro) ──────────────────────────────────────────
      \node[core] (core) {
      \textbf{Core: infrastruttura immutabile}\\[1.5ex]
      {\footnotesize \texttt{Engine} \textbar{} \texttt{DAGScheduler} \textbar{} \texttt{EvidenceStore}}\\[0.8ex]
      {\footnotesize \texttt{TestContext} \textbar{} \texttt{TargetContext}}\\[0.8ex]
      {\footnotesize \texttt{BaseTest} \textbar{} \texttt{BaseConnector} \textbar{} \texttt{BaseGatewayAdapter}}\\[0.8ex]
      {\footnotesize \texttt{TestResult} \textbar{} \texttt{Finding} \textbar{} \texttt{EvidenceRecord}}
      };

      %% ── Moduli Satelliti (A stella) ───────────────────────────────────────
      % Modulo attivo (Sopra)
      \node[active, above=of core] (rest) {
      \textbf{Plugin REST (\gls{OAS})} \ {\footnotesize\itshape (attivo)}\\[1.5ex]
      {\footnotesize \texttt{OpenAPI discovery} \textbar{} \texttt{AttackSurface}}\\[0.8ex]
      {\footnotesize \texttt{SecurityClient} \textbar{} \texttt{KongGatewayAdapter}}\\[0.8ex]
      {\footnotesize 18 Test nativi (Domini D0--D7)}
      };

      % Moduli pianificati (Destra, Sotto, Sinistra)
      \node[planned, right=of core] (gql) {\textbf{GraphQL}\\[0.5ex]{\footnotesize\itshape (pianificato)}};
      \node[planned, below=of core] (grpc) {\textbf{gRPC}\\[0.5ex]{\footnotesize\itshape (pianificato)}};
      \node[planned, left=of core] (ws) {\textbf{WebSocket}\\[0.5ex]{\footnotesize\itshape (pianificato)}};

      %% ── Frecce di Collegamento (Convergenti verso il Core) ────────────────
      \draw[arr-active] (rest.south) -- (core.north);
      \draw[arr-planned] (gql.west) -- (core.east);
      \draw[arr-planned] (grpc.north) -- (core.south);
      \draw[arr-planned] (ws.east) -- (core.west);

      %% ── Legenda e Note ────────────────────────────────────────────────────
      \node[font=\footnotesize, text=gray!80!black, below=4ex of grpc, align=center] (note) {
      frecce continue: modulo attivo \quad frecce tratteggiate: moduli pianificati\\[1ex]
      {\itshape Il discovery dinamico via \texttt{pkgutil} permette il caricamento trasparente di moduli di terze parti}
      };

    \end{tikzpicture}
  \end{adjustbox}
  \caption{Evoluzione di APIGuard Assurance verso una piattaforma modulare: separazione netta tra il core invariante e i plugin di protocollo.}
  \label{fig:modular-platform}
\end{figure}